% ID: SEL-D0034-OB-E04
% ESTADO: validada. S035, cotejada con examen y criterios.
\textbf{a)} La fracción común de muestreo es $375/7500=0.05$.
Multiplicamos por ella la población de cada estrato. En el
orden dado, se seleccionan $\boxed{375,\ 420,\ 285,\ 150}$
personas. En total son $\boxed{1230}$.

\textbf{b)} Tomamos como convención muestras sin reemplazo
y sin orden. Las tres muestras equiprobables son $(2,4)$,
$(2,6)$ y $(4,6)$, con medias 3, 4 y 5. La esperanza es 4 y
la varianza es
\[
\operatorname{Var}(\bar X)=\frac{(3-4)^2+(4-4)^2+(5-4)^2}{3}
=\boxed{2/3}.
\]
Si se considera el orden sin reemplazo, cada media aparece
dos veces y la varianza sigue siendo $2/3$. Si se usa
muestreo con reemplazo, hay nueve pares ordenados equiprobables:
las medias 2, 3, 4, 5 y 6 aparecen 1, 2, 3, 2 y 1 veces.
En ese caso la varianza sería $(4+2+0+2+4)/9=\boxed{4/3}$.
